\documentclass[a4paper]{article}
\usepackage{amsfonts}
\usepackage{amsmath}
\usepackage{amssymb}
\usepackage{graphicx}     

\begin{document}
  
\noindent \large Auteur: Eryk Borczuch \\
\noindent \large Studierichting: Informatica\\
\noindent \large Oefeningenreeks 6A nr. 9.2\\

\medskip

\normalsize

van 3 termen, waarvan de som 33 en het product 1287 is\\

Stel de drie termen op als $(x_2 - v \quad x_2 \quad x_2 + v)$:\\

$(x_2 - v) + x_2 + (x_2 + v) = 33$\\

$(x_2 - v) \cdot x_2 \cdot (x_2 + v) = 1287$\\

Uit de eerste vergelijking:\\

$3x_2 = 33 \rightarrow x_2 = 11$\\

Invullen in de tweede vergelijking:\\

$(11 - v) \cdot 11 \cdot (11 + v) = 1287$\\

$11(121 - v^2) = 1287$\\

$121 - v^2 = 117$\\

$v^2 = 4 \rightarrow v = \pm 2$\\

er zijn dus 2 oplossingen:\\

Als $v = 2$: de rij is $(9, 11 , 13)$\\

Als $v = -2$: de rij is $(13 , 11 , 9)$\\


\begin{thebibliography}{999}
%\bibitem{a} Referentie
\end{thebibliography}
\end{document}